\documentclass[11pt,a4paper]{article}
\usepackage[T1]{fontenc}
\usepackage[utf8]{inputenc}
\usepackage[english]{babel}
\usepackage{amsmath,amssymb}
\usepackage[margin=2cm]{geometry}
\setlength{\parindent}{0pt}
\setlength{\emergencystretch}{3em}
\allowdisplaybreaks

\title{About \texttt{false\_limits.tex}}
\author{}
\date{}
\begin{document}
\maketitle
This note explains what the test data in \texttt{false\_limits.tex} is used for. That file is input for the engine in \texttt{main.c}; it is read by the test \texttt{test\_limits.c} in this folder.

\begin{enumerate}
\item \textbf{Purpose.} The file holds 260 false statements. The engine must never prove any of them, and it must call at least 247 of them false: fewer would mean that a change has made it recognise fewer errors than today.

\item \textbf{Where they come from.} 225 are the \texttt{false\_checks} of \texttt{photos/dati\_esercizi/limiti.jsonl}: the printed result with one integer changed by one, or with its sign flipped, each rejected by the independent numerical oracle. One is the printed result that is wrong (zuccher9-2.7, $-\frac{1}{6}$ instead of $-6$). The other 34 are the same kind of alteration for the statements derived by hand from the exercises without a proposed statement: an infinity with the other sign, a value of the wrong case of a parameter, ``does not exist'' said of a limit that exists, and the converse.

\item \textbf{What it exercises.} That a FALSE verdict is given only when the limit has been computed, and the statement disagrees with it: a number certainly different (its sign comes from interval arithmetic), an infinity where the limit is finite, or a limit where the one-sided limits differ. The 13 statements that are not refuted are the altered versions of the statements in \texttt{book\_limits\_beyond.tex}, which the engine does not decide.

\item \textbf{Expected result.} The lines \texttt{PASS: 260 of 260 false statements rejected} and \texttt{PASS: 247 of 260 statements \dots{} are REFUTED}.

\item \textbf{How to run it.} From the top folder of the repository:
\begin{verbatim}
cc -std=c11 -O2 -Wall -Wextra -o /tmp/test_limits \
    tests/limits/test_limits.c && /tmp/test_limits
\end{verbatim}
\end{enumerate}
\end{document}
